\newcommand{\FourierTransform}[1]{\fouriertransforms \left[#1\right]}
\newcommand{\InverseFourierTransform}[1]{\inverse{\fouriertransforms} \left[#1\right]}

% Semidiscrete Fourier transform
\newcommand{\FourierTransformSemidiscrete}[1]{\fouriertransforms_h \left[#1\right]}
\newcommand{\InverseFourierTransformSemidiscrete}[1]{\inverse{\fouriertransforms}_h \left[#1\right]}

\newcommand{\discreteconvolution}[2]{#1 \ast_h #2}
\renewcommand{\tensorschur}[2]{\ensuremath{#1 \odot #2}} % Schur/Hadamard product

% Discrete Fourier transform
\newcommand{\FourierTransformDiscrete}[1]{\fouriertransforms_h^h \left[ #1 \right]}
\newcommand{\InverseFourierTransformDiscrete}[1]{{\inverse{\fouriertransforms}}_h^h \left[ #1 \right]}
\newcommand{\InverseFourierTransformDiscreteX}[1]{{\inverse{\fouriertransforms}}_h^{h,x} \left[ #1 \right]}

\newcommand{\discreteperiodicconvolution}[2]{#1 \ast_h^h #2}

\newcommand{\FourierTransformDiscreteMatrix}{\tensorq{F}_h^h}

\newcommand{\DiagonalMatrix}{\tensorq{D}}

% Fourier series
\newcommand{\FourierSeries}[1]{\fouriertransforms^h \left[#1\right]}
\newcommand{\InverseFourierSeries}[1]{{\inverse{\fouriertransforms}}^h \left[#1\right]}

% Sine transform
\newcommand{\sinetransforms}{\mathcal{S}}
\newcommand{\SineTransform}[1]{\sinetransforms \left[ #1 \right]}
\newcommand{\InverseSineTransform}[1]{\inverse{\sinetransforms} \left[ #1 \right]}

\newcommand{\SineTransformSemidiscrete}[1]{\sinetransforms_h \left[#1\right]}
\newcommand{\InverseSineTransformSemidiscrete}[1]{\inverse{\sinetransforms}_h \left[#1\right]}

\newcommand{\SineTransformDiscrete}[1]{\sinetransforms_h^h \left[ #1 \right]}
\newcommand{\InverseSineTransformDiscrete}[1]{{\inverse{\sinetransforms}}_h^h \left[ #1 \right]}

\newcommand{\SineTransformDiscreteMatrix}{\tensorq{S}_h^h}

% Fourier series
\newcommand{\SineSeries}[1]{\sinetransforms^h \left[#1\right]}
\newcommand{\InverseSineSeries}[1]{{\inverse{\sinetransforms}}^h \left[#1\right]}

\newcommand{\CosineSeries}[1]{\cosinetransforms^h \left[#1\right]}
\newcommand{\InverseCosineSeries}[1]{{\inverse{\cosinetransforms}}^h \left[#1\right]}

% Cosine transform
\newcommand{\cosinetransforms}{\mathcal{C}}
\newcommand{\CosineTransform}[1]{\cosinetransforms \left[ #1 \right]}
\newcommand{\InverseCosineTransform}[1]{\inverse{\cosinetransforms} \left[ #1 \right]}

\newcommand{\CosineTransformSemidiscrete}[1]{\cosinetransforms_h \left[#1\right]}
\newcommand{\InverseCosineTransformSemidiscrete}[1]{\inverse{\cosinetransforms}_h \left[#1\right]}

\newcommand{\CosineTransformDiscrete}[1]{\cosinetransforms_h^h \left[ #1 \right]}
\newcommand{\InverseCosineTransformDiscrete}[1]{{\inverse{\cosinetransforms}}_h^h \left[ #1 \right]}

% Discrete Dirac function,  bandwidth limited interpolant

\newcommand{\Sdiracdelta}{S}

% J matrix -- matrix with ones on the secondary diagonal
\newcommand{\jmatrix}{\tensorq{J}}
